\subsection{Cross-metric comparison on a single protocol}
\label{sec:res_baselines}

Table~\ref{tab:unified} reports every metric on one reference protocol and one
implementation per metric, which Section~\ref{sec:res_cohort} argues is
necessary for the numbers to be comparable at all. Two rows contain no
generated images. \textsc{real probe} is a second, subject-disjoint block drawn
by the same \textsc{balanced} procedure as the reference: the value a metric
should return when nothing is wrong. \textsc{real head} is the block that the
default filename-ordered protocol selects --- also entirely real BraTS slices,
at the same resolution, from the same collection, preprocessed identically.

\paragraph{A real reference can look like a generative model}
The two real rows are the practical summary of this paper. Drawn correctly, a
real block sits at $\widehat{\mmd}^2 = 0.002$ and FID $= 36.0$. Drawn by the
convenient protocol, the same kind of comparison --- real against real ---
reaches $\widehat{\mmd}^2 = 0.067$ and FID $= 64.6$. The DDPM, an actual
generative model, sits at $0.079$ and $70.8$. On FID the gap between a misdrawn
real reference and a real generator is about six points, set against the
$29$-point spread between the two real rows. A reader given only the scalar
$64.6$ has no way to tell that no generative model was involved in producing
it. This holds for every metric in the table, including the two
natural-image baselines, so it is a property of the protocol rather than of the
medical backbone.

\paragraph{Ordering, and where the backbones disagree}
All four metrics order the three generated sets in the same way relative to the
reference, and all place the weak same-modality generator closer than the
cross-modality shift (Table~\ref{tab:unified}). The metrics disagree on the
\emph{far} end of the ordering. The RadioDINO MMD places retinal fundus further
from brain MRI ($0.463$ at \Lb{12}) than chest CT ($0.403$); FID and KID place
CT further ($322.2$) than retinal ($248.4$). Neither ordering is wrong in the
abstract --- they encode different notions of similarity. RadioDINO, trained on
radiology, treats CT and MRI as neighbouring modalities and ophthalmic
photography as remote; InceptionV3, trained on natural images, responds to
surface appearance, and a colour fundus photograph is closer to its training
distribution than a greyscale CT slice. For evaluating a CT or MRI generator we
regard the radiological ordering as the more useful of the two, but this is an
argument about which similarity a practitioner wants, not evidence that one
backbone is more accurate.

We had previously reported, on a \textsc{head} reference, that CMMD inverts the
ordering of a weak same-modality generator against a cross-modality CT shift.
That measurement was made on the reference protocol this paper argues against,
and we therefore do not carry the claim forward; whether it survives a
subject-balanced reference is left open.

\begin{table}[t]
\centering
\caption{All metrics on one reference protocol (\textsc{balanced}$(100)$,
$N = 500$) and one implementation each. The first two rows contain only real
images. Bootstrap intervals ($60$ resamples) are given for comparisons under
the alternative; none is given for the real rows, where the statistic
approaches degeneracy and the ordinary bootstrap is not consistent
(Section~\ref{sec:cohortnull}).}
\label{tab:unified}
\small
\begin{tabular}{llcccc}
\toprule
Comparison set & Content & $\widehat{\mmd}^2$ \Lb{8} & $\widehat{\mmd}^2$ \Lb{12}
 & FID & KID \\
\midrule
\textsc{real probe}  & real vs.\ real & $0.0016$ & $0.0015$ & $36.0$  & $0.0009$ \\
\textsc{real head}   & real vs.\ real & $0.0628$ & $0.0674$ & $64.6$  & $0.0299$ \\
\midrule
DDPM           & real vs.\ gen. & $0.1028$ & $0.0792$ & $70.8$  & $0.0380$ \\
WDM-3D (MRI)   & real vs.\ gen. & $0.2155$ & $0.2204$ & $174.8$ & $0.1574$ \\
LIDC CT        & real vs.\ gen. & $0.3886$ & $0.4034$ & $322.2$ & $0.3725$ \\
Retinal        & real vs.\ gen. & $0.5026$ & $0.4632$ & $248.4$ & $0.2851$ \\
\bottomrule
\end{tabular}
\end{table}

\paragraph{A note on the degenerate row}
The bootstrap interval our procedure returns for \textsc{real probe} at \Lb{8}
is $[0.0025, 0.0047]$ for a point estimate of $0.0016$ --- an interval that does
not contain the estimate it is meant to bracket. This is not an implementation
error but the expected behaviour of the ordinary bootstrap on a degenerate
$U$-statistic \citep{arcones1992bootstrap}, and it is precisely why we
characterise null behaviour with the empirical inter-cohort distribution of
Section~\ref{sec:cohortnull} rather than with bootstrap intervals. We report the
pathology rather than suppressing it, because a reader who applies the
bootstrap uncritically to a near-null comparison will meet it.
